% ===================================================================
% 3. Methodology
% ===================================================================
\section{Methodology}
\label{sec:method}

Our pipeline is split into three consecutive phases: 2D Object Detection, RGB-only Baseline Pose Estimation, and RGB-D Global Fusion. Each phase builds on top of the previous one, progressively tackling the challenges of 6D pose estimation.



% ------------------------------------------------------------------
\subsection{2D Object Detection}
\label{sec:method_detection}

The first stage of the pipeline is to separate target objects from the background. We fine-tune a YOLO11n~\cite{ultralytics2024yolo11} model on the LineMod dataset~\cite{hinterstoisser2012model}, starting from pre-trained ImageNet weights.

\paragraph{Dataset Format Conversion}
LineMod provides bounding box annotations as $[x, y, w, h]$ (top-left corner plus width and height in pixels). YOLO needs normalized center coordinates $[c_x, c_y, w_n, h_n]$, where all values are in $[0,1]$. We convert them as follows:
\begin{equation}
    c_x = \frac{x + w/2}{W}, \quad c_y = \frac{y + h/2}{H}, \quad w_n = \frac{w}{W}, \quad h_n = \frac{h}{H}
    \label{eq:bbox_conv}
\end{equation}
where $W$ and $H$ are the image dimensions (640$\times$480 for LineMod).

\paragraph{Metric Selection}
Instead of relying only on the standard mAP@50 metric, we focus on the stricter mAP@50-95 (mean AP computed over IoU thresholds from 0.5 to 0.95 in steps of 0.05). Getting a tight bounding box matters a lot for the downstream pose estimation: an oversized crop brings in background noise that can mess up the rotation features extracted by the ResNet backbone.

% ------------------------------------------------------------------
\subsection{RGB-only Baseline Pose Estimation}
\label{sec:method_baseline}

\paragraph{Architecture}
The baseline model uses a ResNet-50~\cite{he2016deep} backbone initialized with ImageNet~\cite{deng2009imagenet} (\texttt{IMAGENET1K\_V1}) pre-trained weights. We remove the original 1000-class classification head (replacing it with \texttt{nn.Identity()}) to expose the 2048-dimensional feature vector. Two parallel linear regression heads are attached:

\begin{itemize}
    \item \textbf{Rotation Head} (\texttt{quat\_head}): $\R^{2048} \rightarrow \R^{4}$, producing a quaternion $\bq = (x, y, z, w)$. An L2-normalization step in the forward pass ensures unit length: $\hat{\bq} = \bq / \|\bq\|_2$.
    \item \textbf{Translation Head} (\texttt{tvec\_head}): $\R^{2048} \rightarrow \R^{3}$, predicting the 3D translation vector $\bt = (t_x, t_y, t_z)$ in meters.
\end{itemize}

\paragraph{Input Pipeline}
Each input image is cropped around the object using a square crop centered on the bounding box: $\text{side} = \max(w, h)$. The crop is resized to $224 \times 224$ pixels with bicubic interpolation and normalized using ImageNet mean and standard deviation. During training, we apply \texttt{ColorJitter} (brightness$=$0.2, contrast$=$0.2), \texttt{GaussianBlur} ($p$=0.2), and \texttt{RandomErasing} ($p$=0.25) to simulate lighting changes, sensor blur, and partial occlusions. We deliberately avoid geometric augmentations (like horizontal flip or rotation) because those would change the ground-truth 3D pose.

\paragraph{Loss Function}
Training both heads at the same time requires balancing their gradients. We use a combined loss:
\begin{equation}
    \mathcal{L}_{\text{baseline}} = \mathcal{L}_{\text{rot}}(\bq_{\text{pred}}, \bq_{\text{gt}}) + \lambda \cdot \mathcal{L}_{\text{trans}}(\bt_{\text{pred}}, \bt_{\text{gt}})
    \label{eq:baseline_loss}
\end{equation}
where the rotation loss uses the quaternion inner product:
\begin{equation}
    \mathcal{L}_{\text{rot}} = 1 - |\bq_{\text{pred}} \cdot \bq_{\text{gt}}|
    \label{eq:rot_loss}
\end{equation}
The absolute value handles the quaternion antipodal ambiguity ($\bq \equiv -\bq$). The translation loss is simply Mean Squared Error (MSE):
\begin{equation}
    \mathcal{L}_{\text{trans}} = \text{MSE}(\bt_{\text{pred}}, \bt_{\text{gt}})
    \label{eq:trans_loss}
\end{equation}
We set the balance parameter $\lambda$ to 1.0. Since $\mathcal{L}_{\text{rot}} \in [0, 1]$ and $\mathcal{L}_{\text{trans}}$ works in squared meters (typical values $\sim 0.01$--$0.1~\text{m}^2$ for LineMod objects at 0.3--0.7\,m distance), the two terms are naturally well-balanced without further tuning.

\paragraph{Pinhole Translation Estimation}
As an extra evaluation baseline, we also implement a purely geometric depth estimate using the pinhole camera model. Given a bounding box with pixel size $s = \max(w, h)$ and the known real-world object diameter $d$, the depth is computed as:
\begin{equation}
    Z = \frac{f_x \cdot d}{s}, \quad X = \frac{(u - p_x) \cdot Z}{f_x}, \quad Y = \frac{(v - p_y) \cdot Z}{f_y}
    \label{eq:pinhole}
\end{equation}
where $(u, v)$ is the bounding box center, $(p_x, p_y)$ is the camera principal point, and $(f_x, f_y)$ are the focal lengths. This approximation implicitly assumes the object is spherical, which leads to large errors for elongated or asymmetric shapes.

% ------------------------------------------------------------------
\subsection{RGB-D Global Fusion}
\label{sec:method_fusion}

To get past the inherent depth ambiguity of using RGB alone, we introduce a multi-modal architecture that combines RGB and depth data. It is inspired by DenseFusion~\cite{wang2019densefusion} but implemented as a simpler 2D global fusion. The full block diagram is shown in Figure~\ref{fig:architecture}.

\begin{figure*}[t]
    \centering
    \resizebox{\textwidth}{!}{%
    \begin{tikzpicture}[
    >={Latex[length=2.2mm]},
    font=\footnotesize,
    node distance=4mm and 12mm,
    blk/.style={draw=gray!55, rounded corners=2pt, fill=gray!8, align=center,
                minimum height=7mm, minimum width=17mm, inner sep=3pt},
    geo/.style={blk, fill=orange!14},
    cat/.style={draw=gray!55, circle, fill=blue!10, align=center,
                minimum size=11mm, inner sep=1pt},
    sumn/.style={draw=gray!55, circle, fill=orange!16, align=center,
                minimum size=8mm, inner sep=1pt},
    spl/.style={draw=gray!55, diamond, aspect=1.4, fill=orange!16, align=center,
                inner sep=1pt, minimum size=7mm},
    fl/.style={->, draw=gray!60, line width=0.5pt},
    gfl/.style={->, draw=orange!70!black, line width=0.6pt},
    dlab/.style={font=\scriptsize\itshape, text=gray!60, inner sep=1.5pt},
    glab/.style={font=\scriptsize\itshape, text=orange!60!black, inner sep=1.5pt},
]
    % inputs
    \node[blk] (rgb) {RGB Crop};
    \node[blk, below=of rgb] (dep) {Depth Crop};
    \node[blk, below=of dep] (met) {Metadata\\BBox \& Camera};
    % feature extractors
    \node[blk, right=of rgb] (r50) {ResNet-50};
    \node[blk, right=of dep] (r18) {ResNet-18};
    \node[blk, right=of met] (enc) {MLP Encoder};
    % fusion path
    \node[cat, right=20mm of r18] (cat) {Concat};
    \node[blk, right=14mm of cat] (fus) {Fusion MLP};
    \node[spl, right=of fus] (spl) {Split};
    % heads (rotation on top, translation at the bottom)
    \node[blk, above right=2mm and 12mm of spl] (rh) {Rotation\\Head};
    \node[blk, below right=2mm and 12mm of spl] (th) {Translation\\Head};
    % outputs
    \node[blk, right=12mm of rh] (gs) {Gram--Schmidt\\Orthogonalization};
    \node[blk, right=of gs] (rout) {Orthogonal\\Rotation $\mathbf{R}$};
    \node[sumn, right=14mm of th] (sum) {$+$};
    \node[blk, right=12mm of sum] (tout) {Translation\\$\bt=\ba+\Delta\bt$};
    % geometric anchor
    \node[geo, below=9mm of enc] (anc) {Anchor $\ba$\\median back-projection\\[1pt]{\scriptsize\itshape no learned parameters}};
    % edges
    \draw[fl] (rgb) -- (r50);
    \draw[fl] (dep) -- (r18);
    \draw[fl] (met) -- (enc);
    \draw[fl] (r50.east) -- (cat) node[dlab, pos=0.45, above]{2048D};
    \draw[fl] (r18.east) -- (cat) node[dlab, midway, above]{512D};
    \draw[fl] (enc.east) -- (cat) node[dlab, pos=0.72, below]{64D};
    \draw[fl] (cat) -- (fus) node[dlab, midway, above]{2624D};
    \draw[fl] (fus) -- (spl) node[dlab, midway, above]{256D};
    \draw[fl] (spl.north) to[out=90, in=180] (rh.west);
    \draw[fl] (spl.south) to[out=-90, in=180] (th.west);
    \draw[fl] (rh) -- (gs) node[dlab, midway, above]{6D};
    \draw[fl] (gs) -- (rout) node[dlab, midway, above]{$3{\times}3$};
    \draw[fl] (th) -- (sum) node[glab, midway, above]{$\Delta\bt$};
    \draw[fl] (sum) -- (tout);
    % anchor wiring
    \coordinate (cA) at ($(met.west)+(-6mm,0)$);
    \coordinate (cB) at ($(met.west)+(-11mm,0)$);
    \draw[gfl] (dep.west) -- (dep.west -| cA) |- ([yshift=2.2mm]anc.west);
    \draw[gfl] (met.west) -- (met.west -| cB) |- ([yshift=-2.2mm]anc.west);
    \draw[gfl] (anc.east) -| (sum.south) node[glab, pos=0.35, above]{$\ba$};
    % group frames
    \begin{scope}[on background layer]
        \node[draw=gray!45, dashed, rounded corners=3pt, inner sep=2.6mm,
              fit=(rgb)(met), label={[font=\scriptsize\sffamily, text=gray!70]above:Inputs}] {};
        \node[draw=gray!45, dashed, rounded corners=3pt, inner sep=2.6mm,
              fit=(r50)(enc), label={[font=\scriptsize\sffamily, text=gray!70]above:Feature Extractors}] {};
    \end{scope}
\end{tikzpicture}%
    }
    \caption{Architecture of the proposed RGB-D Global Fusion model. Three independent feature extractors process the RGB crop (ResNet-50, 2048D), the depth crop (ResNet-18 on a channel-replicated depth map, 512D), and an 8D vector of metadata composed of normalized bounding-box coordinates and camera intrinsics (MLP encoder, 64D). The three vectors are concatenated into a 2624D shared representation, projected to a 256D latent space by the Fusion MLP, and decoded by two heads. The translation head does not output a position: it predicts a residual $\Delta\bt$ that is added to the geometric anchor $\ba$ (orange path), the per-axis median of the depth pixels inside the detected box back-projected through $\mathbf{K}$ (Sec.~\ref{sec:method_anchor}); no learned parameters take part in computing it. The rotation head outputs 6 values mapped to a valid $\SO$ matrix via differentiable Gram-Schmidt orthogonalization.}
    \label{fig:architecture}
\end{figure*}

\paragraph{Depth Preprocessing}
Raw depth maps from the LineMod sensor have invalid pixels (depth$=$0) caused by infrared reflections. We convert raw values to meters using the per-image depth scale from the dataset metadata, and clip to a useful range:
\begin{equation}
    d_{\text{clean}} = \text{clip}\left(\frac{d_{\text{raw}} \cdot s_{\text{depth}}}{1000}, \; 0.0, \; 3.0\right) \; [\text{m}]
    \label{eq:depth_clip}
\end{equation}
The 3.0\,m upper bound covers the entire working range of the LineMod dataset (objects at 0.3--0.7\,m, background $<$ 2\,m) while filtering out outlier sensor readings.

\paragraph{Data Augmentation}
To help the model generalize, the RGB branch gets GPU-accelerated photometric augmentation during training through a custom \texttt{GPUAugmentation} module (built with Kornia): a \texttt{ColorJitter} (brightness, contrast, saturation $=0.1$; hue $=0.05$; $p=0.3$) and \texttt{RandomGaussianNoise} ($\sigma=0.02$; $p=0.2$) are applied on-the-fly on GPU tensors, followed by ImageNet normalization. The depth map is left unaugmented to keep the metric accuracy of the sensor readings. As in the RGB-only baseline (Sec.~\ref{sec:method_baseline}), no geometric augmentations are used.

\paragraph{Global Fusion Architecture}
Our architecture extracts global features from three input modalities through independent branches, as shown in Figure~\ref{fig:architecture}:

\paragraph{RGB Branch}
The first feature extractor is a ResNet-50~\cite{he2016deep} backbone initialized with \texttt{IMAGENET1K\_V1} weights and fine-tuned end-to-end.
The $224\!\times\!224\!\times\!3$ RGB crop goes through the full residual network; the final classification layer is replaced by \texttt{nn.Identity()}, so the global average pooling layer outputs a compact 2048D feature vector $\mathbf{f}_{\text{rgb}} \in \R^{2048}$.
We chose ResNet-50 over shallower alternatives for two reasons.
First, the 2048D bottleneck provides enough capacity to capture the fine-grained texture and appearance details (surface markings, color gradients) needed to distinguish the canonical orientation of textured LineMod objects.
Second, ImageNet pre-training provides generic mid-level features (edges, blobs, textures) that transfer well to the limited diversity of the LineMod training sequences, significantly reducing the risk of overfitting compared to training from scratch.

\paragraph{Depth Branch}
The second extractor processes the cropped and pre-processed depth map.
Since standard pre-trained CNNs expect a 3-channel input, we replicate the single-channel depth map across all three channels before feeding it through a ResNet-18~\cite{he2016deep} backbone, which produces a 512D global average pooling vector $\mathbf{f}_{\text{depth}} \in \R^{512}$.
This channel-replication trick lets us reuse ImageNet weights even for a completely different modality, taking advantage of pre-trained low-level filters (edge detectors) that naturally pick up depth discontinuities and surface normal information encoded as intensity gradients in the depth image.
Compared to a point-cloud-based branch (e.g., PointNet as in DenseFusion), a 2D CNN on the aligned depth crop is cheaper to run and keeps the same architecture style as the RGB branch, simplifying the overall pipeline.


\paragraph{Meta-Branch}
The third branch addresses a subtle but important limitation of purely appearance-based pose estimation: the size and position of an object in the image encode strong geometric hints about its absolute depth and location, but these are hard to extract from convolutional features alone.
To make this information explicit, we build an 8D metadata vector:
\begin{equation}
\mathbf{m} = \left[\frac{c_x}{W},\; \frac{c_y}{H},\; \frac{w}{W},\; \frac{h}{H},\; \frac{f_x}{1000},\; \frac{f_y}{1000},\; \frac{p_x}{W},\; \frac{p_y}{H}\right]
\label{eq:meta}
\end{equation}
where $(c_x, c_y)$ and $(w, h)$ are the bounding-box center and dimensions normalized by the image size $(W, H)$, and $(f_x, f_y, p_x, p_y)$ are the focal lengths and principal point from the camera intrinsic matrix $\mathbf{K}$, also normalized to keep them in a comparable numerical range.
A two-layer MLP ($8 \rightarrow 128 \rightarrow 64$, with ReLU activations) maps $\mathbf{m}$ to a 64D latent vector $\mathbf{f}_{\text{meta}} \in \R^{64}$.
Normalizing each component keeps the MLP inputs in $[0,1]$, preventing gradient imbalances caused by the different physical scales of pixel coordinates and focal-length values.
By feeding camera intrinsics directly into the network, the model can implicitly learn geometry-aware decoding: for example, a larger bounding box combined with a longer focal length clearly signals a smaller depth, a constraint that the convolutional branches alone cannot fully exploit without explicit access to $\mathbf{K}$.

The three feature vectors are concatenated into a single 2624D representation:
\begin{equation}
    \mathbf{f}_{\text{fused}} = [\mathbf{f}_{\text{rgb}} \,\|\, \mathbf{f}_{\text{depth}} \,\|\, \mathbf{f}_{\text{meta}}] \in \R^{2624}
    \label{eq:concat}
\end{equation}
This fused vector is then processed by an MLP ($2624 \rightarrow 512 \rightarrow 256$) to project it into a shared 256D latent space, which feeds the final translation and rotation heads.

\paragraph{Comparison with DenseFusion}
Table~\ref{tab:densefusion_comparison} summarizes the main differences between the original DenseFusion~\cite{wang2019densefusion} and our Global Fusion approach.

\begin{table}[t]
    \centering
    \caption{Comparison between DenseFusion and our Global Fusion architecture.}
    \label{tab:densefusion_comparison}
    \small
    \resizebox{\columnwidth}{!}{%
    \begin{tabular}{@{}lll@{}}
        \toprule
        \textbf{Aspect} & \textbf{DenseFusion} & \textbf{Ours} \\
        \midrule
        Fusion strategy & Dense per-pixel & Global concatenation \\
        Depth processing & PointNet (3D) & CNN on 2D depth map \\
        Rotation repr. & Quaternion & 6D continuous \\
        Iterative refin. & Yes & No \\
        Meta-branch & No & Yes (bbox + $\mathbf{K}$) \\
        \bottomrule
    \end{tabular}}
\end{table}


% ------------------------------------------------------------------
\subsection{Anchored Residual Translation Regression}
\label{sec:method_anchor}

Directly predicting the absolute translation $\bt \in \R^3$ from globally pooled features requires the network to reconstruct a metric position roughly one meter from the camera out of a 256D latent vector --- a task that, as our ablation shows (Sec.~\ref{sec:exp_ablation}), fails when training data is limited under the official LineMod protocol, regardless of the input modalities used. The depth crop, however, already \emph{contains} most of the answer: the visible surface of the object is within a few centimeters of its center.

So we anchor the regression to the observed geometry. Given the detected bounding box and the depth map, we back-project all valid depth pixels inside the (tight) box through the camera intrinsics $\mathbf{K}$ and take the per-axis median:
\begin{equation}
    \ba = \operatorname{median}_{(u,v)\,\in\,\text{box},\ d(u,v)>0}
    \left[ \tfrac{(u-p_x)\,d}{f_x},\ \tfrac{(v-p_y)\,d}{f_y},\ d \right]^\top
    \label{eq:anchor}
\end{equation}
The anchor $\ba$ is a 3D point on the visible surface of the object: on the official training set, its median distance from the ground-truth object center is 19.6\,mm (95th percentile: 49.9\,mm). The translation head then only needs to predict the residual $\Delta\bt$, and the final estimate is $\bt = \ba + \Delta\bt$. This shifts the learning problem from reconstructing an absolute meter-scale position to estimating a centimeter-scale surface-to-center offset --- something that depends only on the object shape and viewpoint, and is therefore learnable from few examples. The median is robust both to background pixels included in the box and to the localization jitter of a real detector; the loss (Sec.~\ref{sec:method_add_loss}) is computed on the composed $\bt$, so training remains end-to-end. If the box contains fewer than 20 valid depth pixels, the anchor falls back to zero and the head gracefully degrades to absolute regression.

% ------------------------------------------------------------------
\subsection{Geometry Input Representations}
\label{sec:method_geometry_inputs}

The depth branch can take the observed geometry in different formats, which we compare in a controlled ablation (Sec.~\ref{sec:exp_ablation}) while keeping the architecture, the number of parameters, and the training recipe identical:

\begin{itemize}
    \item \textbf{Raw depth} (\texttt{raw}): the metric depth crop in meters, channel-replicated to three channels; the translation head predicts the absolute position (no anchor). This is the direct-regression baseline.
    \item \textbf{Anchored normalized depth} (\texttt{norm}): the depth crop centered on the anchor depth and scaled, $d' = \operatorname{clip}\!\big((d - a_z)/s,\,[-4,4]\big)$ with $s = 0.15$\,m ($\approx$ object radius), invalid pixels set to $0$; the head predicts the residual $\Delta\bt$.
\end{itemize}

Depth maps are resized with nearest-neighbor interpolation: interpolating metric values across object boundaries would create 3D points that don't exist anywhere in the actual scene.

% ------------------------------------------------------------------
\subsection{6D Continuous Rotation Representation}
\label{sec:method_rotation6d}

Instead of predicting quaternions or a raw 9D matrix, the rotation head of the fusion model outputs 6 values that are converted to a valid rotation matrix through a differentiable Gram-Schmidt process~\cite{zhou2019continuity}. Given the raw output $\mathbf{a}_1, \mathbf{a}_2 \in \R^3$:
\begin{align}
    \mathbf{b}_1 &= \frac{\mathbf{a}_1}{\|\mathbf{a}_1\|} \label{eq:gs1} \\[3pt]
    \mathbf{b}_2 &= \frac{\mathbf{a}_2 - (\mathbf{b}_1 \cdot \mathbf{a}_2)\,\mathbf{b}_1}{\|\mathbf{a}_2 - (\mathbf{b}_1 \cdot \mathbf{a}_2)\,\mathbf{b}_1\|} \label{eq:gs2} \\[3pt]
    \mathbf{b}_3 &= \mathbf{b}_1 \times \mathbf{b}_2 \label{eq:gs3}
\end{align}
The resulting matrix $\bR = [\mathbf{b}_1 \mid \mathbf{b}_2 \mid \mathbf{b}_3]$ is guaranteed to be in $\SO$ by construction ($\det(\bR) = +1$, orthonormal columns). This representation is continuous and differentiable, which means stable gradient flow during backpropagation.

% ------------------------------------------------------------------
\subsection{Loss Function for RGB-D Fusion}
\label{sec:method_add_loss}

Unlike the combined loss of the baseline (Eq.~\ref{eq:baseline_loss}), the fusion model is trained directly with the ADD (Average Distance of Model Points) metric as the loss:
\begin{equation}
    \mathcal{L}_{\text{ADD}} = \frac{1}{N} \sum_{i=1}^{N} \|(\bR_{\text{pred}} \cdot \bp_i + \bt_{\text{pred}}) - (\bR_{\text{gt}} \cdot \bp_i + \bt_{\text{gt}})\|_2
    \label{eq:add_loss}
\end{equation}
where $\{\bp_i\}_{i=1}^N$ are $N$=500 points sampled from the 3D CAD model. This directly optimizes the same metric used for evaluation, removing the need for the ad-hoc balance parameter $\lambda$ between rotation and translation losses. We do not use this loss for the RGB-only baseline because its inherent depth estimation errors would cause huge metric distances (on the order of meters), leading to exploding gradients and unstable training. Its combined loss, by isolating the rotation error (bounded at 1.0), serves as a key stabilization mechanism.

